%% Plantilla de minuta de reunión
%% Uso: copiar como minuta-NNN-<tema>.tex, completar y compilar.
%% Compilar: tectonic -X compile plantilla-minuta.tex
\documentclass[11pt,a4paper]{article}
\usepackage{../latex/legasoft}

\lgtitulo{Minuta de reunión}
\lgsubtitulo{\campo{Asunto de la reunión}}
\lgtituloCorto{Minuta \campo{NNN}}
\lgcodigo{LGS-MIN-NNN}
\lgversion{0.1}
\lgestado{Plantilla}
\lgfecha{\campo{dd-mm-aaaa}}
\lgautor{\lgDirector\ (\lgDirectorRol)}

\begin{document}

% Una minuta es un documento corto: no lleva portada ni índice.
\begin{center}
{\sffamily\bfseries\LARGE\color{lgPrimario} Minuta de reunión\par}
\vspace{2pt}
{\sffamily\normalsize\color{lgGris} Legasoft · \campo{asunto}\par}
\end{center}
\vspace{6pt}
{\color{lgBorde}\hrule height 0.8pt}
\vspace{10pt}

\begin{porcompletar}[Cómo usar esta plantilla]
Copie este archivo como \ruta{minuta-NNN-<tema>.tex}, sustituya los campos en
ámbar y elimine esta caja antes de distribuir la minuta. La preparación y
conservación de las minutas corresponde al \lgDirectorRol\ según
\texttt{LGS-INST-001}, sección 10.1.
\end{porcompletar}

% =============================================================================
\section*{1. Datos de la reunión}

\begin{tabularx}{\linewidth}{@{}L{3.6cm} Y@{}}
\toprule
\sffamily\bfseries\small\color{lgPrimario}Código        & LGS-MIN-\campo{NNN} \\
\sffamily\bfseries\small\color{lgPrimario}Fecha y hora  & \campo{dd-mm-aaaa}, de \campo{hh:mm} a \campo{hh:mm} \\
\sffamily\bfseries\small\color{lgPrimario}Modalidad     & \campo{presencial / virtual / mixta} \\
\sffamily\bfseries\small\color{lgPrimario}Lugar o enlace& \campo{sala o plataforma} \\
\sffamily\bfseries\small\color{lgPrimario}Tipo          & \campo{interna del equipo / con la parte interesada} \\
\sffamily\bfseries\small\color{lgPrimario}Facilitador   & \lgDirector \\
\sffamily\bfseries\small\color{lgPrimario}Redactor      & \campo{nombre} \\
\bottomrule
\end{tabularx}

% =============================================================================
\section*{2. Participantes}

\begin{tabularx}{\linewidth}{@{}L{4.6cm} Y C{2.2cm}@{}}
\toprule
\sffamily\bfseries\color{lgPrimario}Nombre &
\sffamily\bfseries\color{lgPrimario}Rol u organización &
\sffamily\bfseries\color{lgPrimario}Asistencia \\
\midrule
\lgDirector   & \lgDirectorRol   & \campo{sí / no} \\
\lgArquitecto & \lgArquitectoRol & \campo{sí / no} \\
\lgQA         & \lgQARol         & \campo{sí / no} \\
\campo{nombre} & \campo{organización cliente} & \campo{sí / no} \\
\bottomrule
\end{tabularx}

% =============================================================================
\section*{3. Agenda}

\begin{enumerate}
  \item \campo{Punto de agenda}
  \item \campo{Punto de agenda}
  \item \campo{Punto de agenda}
\end{enumerate}

% =============================================================================
\section*{4. Desarrollo de la reunión}

Resumen de lo tratado en cada punto, con la información relevante para
reconstruir el contexto de las decisiones.

\begin{porcompletar}
\campo{Punto 1: qué se expuso, qué se discutió y qué se concluyó.}

\campo{Punto 2: \ldots}
\end{porcompletar}

% =============================================================================
\section*{5. Acuerdos y decisiones}

\begin{tabularx}{\linewidth}{@{}C{1.8cm} Y L{3.2cm}@{}}
\toprule
\sffamily\bfseries\color{lgPrimario}ID &
\sffamily\bfseries\color{lgPrimario}Acuerdo o decisión &
\sffamily\bfseries\color{lgPrimario}Justificación \\
\midrule
A-01 & \campo{acuerdo} & \campo{motivo} \\
A-02 & \campo{acuerdo} & \campo{motivo} \\
\bottomrule
\end{tabularx}

% =============================================================================
\section*{6. Compromisos y seguimiento}

Cada compromiso tiene un responsable y una fecha; su seguimiento se registra en
Jira, conforme a la responsabilidad de coordinación del proyecto.

\begin{tabularx}{\linewidth}{@{}C{1.8cm} Y L{3.6cm} C{2.2cm} C{2.0cm}@{}}
\toprule
\sffamily\bfseries\color{lgPrimario}ID &
\sffamily\bfseries\color{lgPrimario}Compromiso &
\sffamily\bfseries\color{lgPrimario}Responsable &
\sffamily\bfseries\color{lgPrimario}Fecha límite &
\sffamily\bfseries\color{lgPrimario}Estado \\
\midrule
C-01 & \campo{acción concreta} & \campo{nombre} & \campo{dd-mm-aaaa} & \campo{pendiente} \\
C-02 & \campo{acción concreta} & \campo{nombre} & \campo{dd-mm-aaaa} & \campo{pendiente} \\
\bottomrule
\end{tabularx}

% =============================================================================
\section*{7. Temas pendientes y riesgos detectados}

\begin{porcompletar}
\campo{Asuntos que quedaron sin resolver, restricciones nuevas o riesgos que
deban incorporarse al registro del proyecto.}
\end{porcompletar}

% =============================================================================
\section*{8. Próxima reunión}

\begin{tabularx}{\linewidth}{@{}L{3.6cm} Y@{}}
\toprule
\sffamily\bfseries\small\color{lgPrimario}Fecha prevista & \campo{dd-mm-aaaa} \\
\sffamily\bfseries\small\color{lgPrimario}Objetivo       & \campo{propósito de la siguiente sesión} \\
\bottomrule
\end{tabularx}

\vspace{1.4cm}
\begin{center}
\begin{tabularx}{\linewidth}{@{}Y Y@{}}
\centering\rule{5.4cm}{0.5pt} &
\centering\arraybackslash\rule{5.4cm}{0.5pt} \\[-2pt]
\centering\sffamily\small\lgDirector &
\centering\arraybackslash\sffamily\small\campo{representante del cliente} \\[-4pt]
\centering\scriptsize\color{lgGris}Por Legasoft &
\centering\arraybackslash\scriptsize\color{lgGris}Conformidad con lo registrado \\
\end{tabularx}
\end{center}

\end{document}
